\section{Truncated Hermite hierarchies}\label{sec:truncated}

Fourier--Hermite codes for KRMHD evolve the moments $g_0,\dots,g_M$ of a branch and close the hierarchy at $m=M$; examples are Viriato \citep{Loureiro2016} and GANDALF \citep{Kanekar2025}. The truncated hierarchy consists of exactly $M+1$ passive slices, whose phase velocities and weights, for the zero closure, are the nodes and weights of the Gauss rule of the spectral density $\varrho\,\dd u$ of (\ref{eq:varrho}). Appendix~\ref{app:truncated} gives the proofs.

\subsection{The slices of the truncated hierarchy}

For one branch truncated at $m=M$, the vector $g=(g_0,\dots,g_M)^{\mathsf T}$ obeys (\ref{mod:hierarchy}), where, for the zero closure $g_{M+1}=0$, $A$ is the leading $(M+1)\times(M+1)$ block of the ladder matrix (\ref{mod:ladder}). The copy closure $g_{M+1}=g_{M-1}$ changes, for $M\ge2$, only $A_{M,M-1}$, to $\sqrt{M/2}+\sqrt{(M+1)/2}$. The weight $\Pw$ of section~\ref{sec:model} makes $\Pw A$ symmetric for the zero closure; for the copy closure its last entry becomes $p_M=\sqrt M/(\sqrt M+\sqrt{M+1})\approx\frac12$, and $\Pw A$ is then the same matrix for both closures.

\begin{theorem}\label{thm:trunc}
Let $\Lambda\notin[0,1]$, the value $\Lambda=\infty$ included, and use the zero closure or, for $M\ge2$, the copy closure.
\begin{itemize}
\item[(i)] $A$ has $M+1$ real and simple eigenvalues $u_0<\dots<u_M$. The right eigenvector of $u_j$ is $\boldsymbol r_j=(r_0(u_j),\dots,r_M(u_j))^{\mathsf T}$, where $r_m(u)$, the $m$th Hermite coefficient of the Case--Van Kampen mode $F_u$ of section~\ref{sec:slices}, is a polynomial of degree $m$. The left eigenvector is $\boldsymbol w_j=\Pw\boldsymbol r_j$, and $\boldsymbol w_j^{\mathsf T}\boldsymbol r_l=n_j\,\delta_{jl}$ with $n_j=\boldsymbol r_j^{\mathsf T}\Pw\boldsymbol r_j>0$.
\item[(ii)] For arbitrary real $\Phi$ and $\Psi$, and brackets in (\ref{mod:hierarchy}) evaluated by any bilinear rule, possibly followed by linear operations, such as dealiasing, that act on every $g_m$ in the same way, $g$ solves (\ref{mod:hierarchy}) if and only if the unit-weight slice amplitudes $a_j=\boldsymbol w_j^{\mathsf T}g/\sqrt{n_j}=\sum_mp_mR_{mj}g_m$, with $R_{mj}=r_m(u_j)/\sqrt{n_j}$, satisfy
\begin{equation}\label{eq:truncslice}
  \partial_ta_j+\{\Phi+su_j\Psi,\,a_j\}+su_j\,\partial_za_j=0,\qquad j=0,\dots,M,
\end{equation}
with the brackets evaluated by the same rule.
\item[(iii)] $g_m=\sum_jR_{mj}\,a_j$, and $\sum_mp_m|g_m(\boldsymbol k)|^2=\sum_j|a_j(\boldsymbol k)|^2$ for every Fourier mode.
\end{itemize}
\end{theorem}

Each slice is a passive scalar carried by $\uE+su_j\bhat$, the velocity $\boldsymbol U_u$ of (\ref{eq:Uu}) at $u=u_j$, so theorem~\ref{thm:trunc} is the truncated form of theorem~\ref{thm:slice}, for any $z^\pm$. Its proof uses only the bilinearity of the discrete bracket. The conservation of $\int a_j^2\,\dd^3r$ and the two-slice law (\ref{eq:twoslice}) need, in addition, the antisymmetry of the bracket under the volume average, which the dealiased pseudo-spectral bracket has exactly only when $\Phi$, $\Psi$ and $g$ lie inside the dealiasing band.

\citet[\S6.2.1]{Loureiro2016} diagonalised the $k_\perp$-dependent truncated streaming matrix of the KREHM model for the linear parallel step of Viriato. \citet{ParkerDellar2015} identified the eigenvectors of the truncated Hermite system of linearised Vlasov--Poisson with discrete Case--Van Kampen modes. The normal fields of truncated Hermite hierarchies of electromagnetic drift-kinetic models \citep{Tassi2015} are, as \citet[\S7]{Tassi2022} showed, proportional to the truncated distribution function at the zeros of a Hermite polynomial and, in three dimensions, coupled through the parallel gradients of the self-consistent potentials. In KRMHD the Landau coupling enters only $A_{10}$ of a matrix that does not depend on the wavenumber, and the Alfv\'enic fields do not depend on $g$, so the decoupling is exact in three dimensions.

\subsection{Gauss quadrature of the Case--Van Kampen measure}

\begin{proposition}\label{tr:gaussprop}
Let $\Lambda\notin[0,1]$ and use the zero closure. Then (i) $\int\varrho\,r_mr_n\,\dd u=\delta_{mn}/p_m$; (ii) the $u_j$ are the zeros of $r_{M+1}$, and with $\lambda_j=1/n_j$
\begin{equation}\label{eq:gauss}
  \sum_{j=0}^M\lambda_j\,f(u_j)=\int\varrho(u)\,f(u)\,\dd u\qquad\text{for every polynomial $f$ of degree at most $2M+1$,}
\end{equation}
so the $u_j$ and $\lambda_j$ are the nodes and weights of the $(M+1)$-point Gauss rule of $\varrho\,\dd u$; (iii) for $G=\sum_{m\le M}g_me_m$, with $e_m=H_mF_M/\sqrt{2^mm!}$, $a_j=\sqrt{\lambda_j}\,\Tr[G](u_j)$.
\end{proposition}

Since $\lambda_j\simeq\varrho(u_j)\Delta_j$, with $\Delta_j$ the local spacing of the nodes, part (iii) gives $a_j\simeq\Delta_j^{1/2}a_{u_j}$ for states in the span of the retained moments: the discrete slices are the continuum slices (\ref{eq:aunit}) at the nodes. For smooth states outside that span, the distribution of free energy over the slices converges to the continuum one approximately as $M^{-1}$.

At $\Lambda=\infty$ this is the known correspondence with Gauss--Hermite quadrature: the nodes $x_j$ are the zeros of $H_{M+1}$, as \citet{ParkerDellar2015} and \citet{Tassi2022} noted, with weights $\lambda^\infty_j$, and the link with the Jacobi matrix is that of \citet{GolubWelsch1969}. The identification of the truncated KRMHD hierarchy at finite $\Lambda$ with the Gauss rule of the dressed density $\varrho=F_M/|\eps|^2$ is new, as is the exact nonlinear decoupling. The copy closure gives no Gauss rule: its nodes are the roots of $r_{M+1}=r_{M-1}$ (\ref{appC:eigcond}), not of $r_{M+1}$, and the rule with weights $1/n_j$ is exact only to degree $2M-1$.

\subsection{Conserved quadratic forms}

The truncated counterpart of (\ref{eq:Qphi}) is $Q^{(M)}_\varphi=\int g^{\mathsf T}\Pw\varphi(A)g\,\dd^3r=\sum_j\varphi(u_j)\int a_j^2\,\dd^3r$.

\begin{proposition}
Let $\Lambda\notin[0,1]$, use either closure, let $\Phi$, $\Psi$ and $g$ lie inside the dealiasing band, and let no collisions, dissipation or forcing act on $g$. A form $\int g^{\mathsf T}Kg\,\dd^3r$ with a real symmetric matrix $K$ is conserved for all $\Phi$ and $\Psi$ if and only if $K=\Pw\varphi(A)$ for a real function $\varphi$ on the nodes.
\end{proposition}

\begin{proof}
By the antisymmetry of the bracket and of $\gradpar$, the form changes only through $B=\frac12(KA-A^{\mathsf T}K)$, and $B_{mn}\neq0$ would make the state $g_m=\cos k_zz$, $g_n=\sin k_zz$, $\Psi=0$ change it. So $KA=A^{\mathsf T}K$, $R^{\mathsf T}KR$ commutes with $\mathrm{diag}(u_j)$, whose entries are distinct, and $K=\Pw R\,\mathrm{diag}(\varphi_j)R^{\mathsf T}\Pw=\Pw\varphi(A)$. Conversely, each $\int a_j^2\,\dd^3r$ is conserved by (\ref{eq:truncslice}).
\end{proof}

For the zero closure and $G=\sum_{m\le M}g_me_m$, $\fe{G}{G}=g^{\mathsf T}\Pw g$ and $\fe{G}{LG}=g^{\mathsf T}\Pw Ag$, so by (\ref{eq:WGH}) $W^\pm=(n_iT_i/2V)\,Q^{(M)}_1$ and $\Gamma^\pm=(\vth/2V)\,Q^{(M)}_u$ are exactly CMM's (B31) and (B32) on the retained moments. Under the copy closure $Q^{(M)}_1$ carries the weight $p_M$ on $g_M^2$ and is not (B31); $Q^{(M)}_u$ is the same for both closures.

\subsection{The phase-space helicity at odd and even \texorpdfstring{$M$}{M}}

Let $h_m=H_m/\sqrt{2^mm!}$, $\psi_m=h_mF_M^{1/2}$, $J$ the ladder matrix of the zero closure at $\Lambda=\infty$, and $\boldsymbol h(0)$ and $\boldsymbol h'(0)$ the vectors with components $h_m(0)$ and $h_m'(0)$. For the zero closure, the exactness of Gauss--Hermite quadrature gives the matrix $(B_M)_{mn}=\PVint\psi_m\psi_nu^{-1}\,\dd u$ of CMM's helicity (\ref{eq:WGH}) on the retained moments,
\begin{equation}\label{eq:Hoddeven}
  B_M=J^{-1}=\Pw A^{-1}\quad(M\ \text{odd}),\qquad
  B_M=J^{+}+\lambda^\infty_{M/2}\big[\boldsymbol h(0)\,\boldsymbol h'(0)^{\mathsf T}+\boldsymbol h'(0)\,\boldsymbol h(0)^{\mathsf T}\big]\quad(M\ \text{even}),
\end{equation}
with $J^+$ the Moore--Penrose inverse. For odd $M$ the conserved form $Q^{(M)}_{1/u}$ is therefore exactly CMM's helicity of the truncated state, $\Hph^\pm=(n_iT_i/2\vth V)\,Q^{(M)}_{1/u}$, for every $\Lambda\neq1$.

For even $M$, $u_{M/2}=0$ is an eigenvalue for every $\Lambda$ and both closures, because $SAS=-A$ with $S=\mathrm{diag}((-1)^m)$ and $M+1$ is odd, and for the zero closure the slice at $u=0$ is proportional to the truncated distribution at $v_\parallel=0$ (appendix~\ref{app:truncated}). It is advected by $\uE$ alone, and its conserved energy is proportional to the imaginary part of (\ref{eq:Hcomplex}) for the truncated state, the truncated form of CMM's conservation of $\mathrm{Im}\,\Hph$ (their eqs.~2.22--2.23). The real part is not conserved: the rank-two term in (\ref{eq:Hoddeven}) couples the null vector $\boldsymbol h(0)$ of $J$ to its other eigenvectors, so $B_M$ is not of the form $\Pw\varphi(A)$ (at finite $\Lambda$ we found this numerically). The conserved principal-value form $\sum_{u_j\neq0}u_j^{-1}\int a_j^2\,\dd^3r$ leaves out the slice at $u=0$. At $\Lambda=\infty$ it differs from $\int g^{\mathsf T}B_Mg\,\dd^3r$ by exactly the node term $-2\lambda^\infty_{M/2}\int[\boldsymbol h(0)\cdot g][\boldsymbol h'(0)\cdot g]\,\dd^3r$ (appendix~\ref{app:truncated}), with $\lambda^\infty_{M/2}\simeq[\pi/(2M+3)]^{1/2}$. At finite $\Lambda$ the difference is another rank-two form of the same order, and in both cases the even-$M$ form converges to the helicity only as $M^{-1/2}$; for $G=F_M(u-\frac12)$ the difference is between 5\% and 18\% of the helicity at $M=64$ and 128. Diagnostics of the helicity should therefore use odd $M$ and the zero closure.

\citet{Tassi2022} found the field advected by the $\boldsymbol E\times\boldsymbol B$ flow alone at an odd number of retained moments, and its Casimir invariants in three dimensions, for Hermite truncations of drift-kinetic models with a symmetric matrix; \citet[\S3.4]{Tassi2023} restates the field for gyrofluid reductions, and \citet{Tassi2015} gives it in two dimensions. New here are the survival of this field under the $\Lambda$-modified ladder and both closures, its identification with $\mathrm{Im}\,\Hph$, and the consequence (\ref{eq:Hoddeven}) for the real part.

\subsection{The largest phase velocity and the time step}

The largest phase velocity $u_{\max}=u_M$ hardly depends on $\Lambda$. At $\Lambda=\infty$ it is the largest Gauss--Hermite node $x_M$, and since $\Lambda$ acts only on the moments $m=0$ and 1, where the extreme eigenvector of $J$ is small, the shift at finite $\Lambda$ is of size $x_Me^{-x_M^2}/\Lambda$ up to an algebraic factor, about $10^{-49}$ at $M=64$. With $\kappa_m=(2m+1)^{1/2}$,
\begin{equation}\label{eq:umax}
\begin{aligned}
  u_{\max}&=\kappa+2^{-1/3}\alpha_1\,\kappa^{-1/3}-\frac{2^{-2/3}\alpha_1^2}{10}\,\kappa^{-5/3}+O(\kappa^{-3}),&\qquad\kappa&=\kappa_{M+1}\quad\text{(zero closure)},\\
  u_{\max}&=\kappa+2^{-1/3}\alpha_1'\,\kappa^{-1/3}-2^{-5/3}\Big(\frac{4}{5\alpha_1'}+\frac{\alpha_1'^2}{5}\Big)\kappa^{-5/3}+O(\kappa^{-3}),&\qquad\kappa&=\kappa_M\quad\text{(copy closure)},
\end{aligned}
\end{equation}
where $\alpha_1=-2.3381074$ and $\alpha_1'=-1.0187930$ are the first zeros of the Airy function and of its derivative. The zero-closure line is the classical expansion of the largest zero of $H_{M+1}$: its first two terms are given by \citet[theorem~6.32]{Szego1975} \citep[see also][eq.~18.16.17]{DLMF}, and the third by \citet{Olver1959} \citep[eqs~12.7.2 and~12.11.9]{DLMF}. The copy-closure line is new; appendix~\ref{app:truncated} sketches its derivation. Figure~\ref{fig:spectrum}(b) compares (\ref{eq:umax}) with the eigenvalues.

By theorem~\ref{thm:trunc} the nonlinear term $-\{\Phi,g\}-s\{\Psi,Ag\}$ advects slice $j$ with the stream function $\Phi+su_j\Psi$, at a speed up to $\max|\nabla_\perp\Phi|+s|u_j|\max|\nabla_\perp\Psi|$. A time-step limit for an explicit treatment of this term must therefore use $su_{\max}|\nabla_\perp\Psi|$, not $s|\nabla_\perp\Psi|$. \citet[eq.~5.3]{Mandell2024} include the largest Gauss--Hermite node in the parallel-streaming frequency that sets the time step of the GX code; here the same factor also multiplies the field-line term. \citet[eq.~2.26]{KanekarThesis2015} had already reduced the nonlinear time step of a Fourier--Hermite KRMHD code by the factor $M^{-1/2}$, which scales as $1/u_{\max}$, without deriving it.
